% 来源：OEBPS/Text/01_17.xhtml（OPF spine 顺序；仅第一册）
\cleardoublepage
\BookEntry[达尔文《物种起源》]{达尔文《物种起源》}

\BookCover{images/书单04物种起源.jpg}{[英] 达尔文著| 周建人等译| 商务印书馆1995年版}

\begin{BookSummaryBox}
{\centering \textbf{19世纪关于自然科学的三大发现之一\\ 被誉为“影响历史进程的书”}\par}
\end{BookSummaryBox}

\subsubsection*{这本书谈什么？}

《物种起源》是达尔文论述生物进化的重要著作，首次出版于1859年11月24日，全称《论借助自然选择（即在生存斗争中保存优良族）的方法的物种起源》。
\par

达尔文1831年从剑桥大学毕业后，以博物学家的身份乘海军勘探船“贝格尔号”作历时5年(1831\textasciitilde{}1836年)的环球旅行，观察和搜集了动物、植物和地质等方面的大量材料，经过归纳整理和综合分析，形成了生物进行的概念。首次提出了进化论的观点，即物种的演化是通过自然选择（天择）和人工选择（人择）的方式实现的。自然选择的进化学说对各种唯心的神造论、目的论和物种不变论提出根本性的挑战。使当时生物学各领域已经形成的概念和观念发生根本性的改变。
\par

\subsubsection*{芒格为什么谈起这本书？}

在南加州大学GOULD法学院毕业典礼上的演讲中，芒格谈到他有一条“铁律”，帮助自己在偏向于某种强烈的意识形态时保持清醒：
\par

\begin{quotation}
我觉得我没资格拥有一种观点，除非我能比我的对手更好地反驳我的立场。
\par
\end{quotation}

达尔文也有自己的“铁律”。达尔文早在19世纪30年代末期就形成了他的物种进化理论，但直到1859年，他才出版了他的巨著《物种起源》。达尔文认为，任何对人类起源提出新解释的科学理论都将会遭遇广泛的歧视，所以他在公开发表他的理论之前，先谨慎地对各种可能出现的反驳进行了准备。因此，他花了整整20年，努力地完善他的理论，并为这种理论做好了辩护的准备。
\par

\subsubsection*{《物种起源》作者达尔文如是说}

“认为生命及其若干能力原来是由‘造物主’注入到少数类型或一个类型中去的，而且认为在这个行星按照引力的既定法则继续运行的时候，最美丽的和最奇异的类型，从如此简单的始端，过去、曾经而且现今还在进化着——这种观点是极其壮丽的。”
\par

“人类在道德文化方面最高级的阶段，就是当我们认识到应当用理智控制思想时。”
\par

“乐观是希望的明灯，它指引着你从危险峡谷中步向坦途，使你得到新的生命和新的希望，支持着你的理想永不泯灭。”
\par

“敢于浪费哪怕一个钟头时间的人，说明他还不懂得珍惜生命的全部价值。”
\par
